\section{Convex dispatch can have negative marginal prices}
\label{app:networkmonotonicity}
The monotonicity premise in Theorem~\ref{thm:lossyfrontier} is substantive
when supply cost is obtained by optimizing a network dispatch. Transmission
congestion can yield negative nodal marginal prices even with nonnegative
generation offers; this is an established effect
\citep[Section V-B, Figure 5]{bakerGangammanavar2024prices}.
The following independently checked symbolic example makes the distinction
explicit with strictly increasing, strictly convex generator costs. It is
a scope illustration, not a new negative-price result or empirical network.

Consider a lossless three-bus DC network with unit susceptance on each
of lines $12,13,23$, oriented flows $f_{ij}=\theta_i-\theta_j$, and
$\theta_3=0$. Symmetric line limits are $5,1/2,5$, respectively.
Background demand is $(4,0,1)$ and an additional depot withdrawal
$x\in[0,1]$ occurs at bus 3. Generators at buses 1 and 2 satisfy
$0\le g_1,g_2\le10$, with costs
\[
 C_1(g_1)=3g_1+g_1^2/100,\qquad
 C_2(g_2)=g_2+g_2^2/100.
\]
All quantities use consistent power and per-period cost units; a fixed
period converts them to energies. Nodal balance is generation minus
demand equal to outgoing flow. Let $J(x)$ be the minimum total generation
cost under these constraints.

Bus-3 balance and total balance give
$\theta_1+\theta_2=1+x$ and $g_2=5+x-g_1$. Bus-1 balance then implies
\[
 \theta_1=(g_1-3+x)/3,\qquad g_1\ge3/2-x,
\]
where the inequality follows from the lower limit on line 13. The point
\[
 g_1=3/2-x,\quad g_2=7/2+2x,\quad
 \theta_1=-1/2,\quad\theta_2=3/2+x
\]
attains that lower bound and is feasible throughout the stated interval:
its flows are $(-2-x,-1/2,3/2+x)$, and all generator limits hold.
After eliminating $g_2$, the objective derivative with respect to $g_1$ is
\[
 2+\frac{2g_1-5-x}{50}\ge2-\frac6{50}>0.
\]
Thus every feasible point has cost at least that of the displayed point;
this establishes global optimality, not merely stationarity at a proposed
active set. Direct expansion yields
\[
 J(x)-J(0)=-\frac{89}{100}x+\frac{x^2}{20}.
\]
This incremental cost is strictly convex and strictly decreasing on
$[0,1]$. More withdrawal at bus 3 relieves the congested reverse flow
toward bus 1, reducing output at the expensive generator at bus 1 while
increasing cheaper generation at bus 2. The redispatch saving exceeds the cost of the
extra supply. Both generators' own cost derivatives remain positive.

Consequently convexity of the optimized supply value and positive
generation marginal costs do not ensure monotonicity in nodal load.
A network application of the efficient-frontier reduction must establish
the required load-domain monotonicity, or retain the alternatives that
coordinatewise dominance would discard. The example does not invalidate
the theorem under its hypotheses, change the public synthetic supply
objective, or establish AC feasibility, prevalence or a learning advantage.
